\documentclass[12pt]{article}
\usepackage[utf8]{inputenc}
\usepackage[T1]{fontenc}
\usepackage[margin=1in]{geometry}
\usepackage{setspace}
\usepackage{xcolor}
\usepackage[colorlinks=true, urlcolor=blue!60!black, linkcolor=blue!60!black]{hyperref}
\setlength{\parindent}{0em}
\setlength{\parskip}{0.45em}
\renewcommand{\baselinestretch}{1.0}
\begin{document}
\begin{center}
{\Large \textit{Crosby v. National Foreign Trade Council}, 530 U.S.\ 363 (2000)\par}
\vspace{0.3em}
{\normalsize Case Brief --- One-Page Edition\par}
\end{center}

\textbf{Memory Jogger.} Massachusetts's Burma boycott was preempted by Congress's calibrated Burma sanctions (\href{https://supreme.justia.com/cases/federal/us/530/363/}{Justia}).

\textbf{Facts.} Massachusetts penalized state bidders ``doing business with Burma'' unless they underbid others by 10\%, reaching foreign firms and existing investments, with no waiver or sunset. Congress then enacted narrower sanctions, presidential waiver authority, and a multilateral strategy. The EU and Japan protested the state law.

\textbf{Procedural History.} The National Foreign Trade Council sued; the district court and First Circuit enjoined the law on foreign affairs, Foreign Commerce Clause, and preemption grounds.

\textbf{Votes.} Affirmed, 9--0: Souter for the Court; Scalia and Thomas concurring in the judgment.

\textbf{Holding.} The state law is an obstacle to the federal Act's presidential discretion, its limits on whom and what sanctions reach, and its diplomacy.

\textbf{Reasoning.} Purchasing power is no defense; the rigid law reduced the President's ``bargaining chips,'' penalized what Congress permitted, and impeded diplomacy. Congressional silence was ``ambiguous,'' not permission. Scalia rejected the opinion's needless use of legislative history.

\textbf{Examples.} \emph{Same-Side:} a broader, unwaivable state ban. \emph{Opposite-Side:} a boycott with no federal statute in the field. \emph{Fence-Sitter:} divestment that tracks a federal list.

\textbf{Critique.} Treating any partial sanction as a ``calibration of force'' invites expansive preemption, and crediting foreign protests risks a foreign veto over state policy.

\textbf{Key Quotation.} ``The fact of a common end hardly neutralizes conflicting means, see \textit{Gade v.\ National Solid Wastes Management Assn.}, 505 U.S.\ 88, 103 (1992), and the fact that some companies may be able to comply with both sets of sanctions does not mean that the state Act is not at odds with achievement of the federal decision about the right degree of pressure to employ.'' --- Justice Souter, 530 U.S.\ at 379.

\end{document}
